
\documentclass{article}
\usepackage{colortbl}
\usepackage{makecell}
\usepackage{multirow}
\usepackage{supertabular}

\begin{document}

\newcounter{utterance}

\centering \large Interaction Transcript for game `function\_detective', experiment `numbers\_passive\_examples\_oneshot', episode 10 with gpt{-}5.4.
\vspace{24pt}

{ \footnotesize  \setcounter{utterance}{1}
\setlength{\tabcolsep}{0pt}
\begin{supertabular}{c@{$\;$}|p{.15\linewidth}@{}p{.15\linewidth}p{.15\linewidth}p{.15\linewidth}p{.15\linewidth}p{.15\linewidth}}
    \# & \multicolumn{2}{c}{Player} && \multicolumn{2}{c}{Game Master} \\
    \hline

    \theutterance \stepcounter{utterance}  
    & & & \multicolumn{4}{p{0.6\linewidth}}{
        \cellcolor[rgb]{0.9,0.9,0.9}{
            \makecell[{{p{\linewidth}}}]{
                \texttt{\tiny{[P1$\langle$GM]}}
                \texttt{I have a hidden Python function for you to figure out.} \\
\\ 
\texttt{Your job is to infer what it does and then submit a correct Python implementation.} \\
\\ 
\texttt{Function signature:} \\
\texttt{(x: int) {-}> int} \\
\\ 
\texttt{You have 1 turns total.} \\
\\ 
\texttt{I have already given you a fixed list of valid input{-}output examples. Solve using those examples only.} \\
\\ 
\texttt{You must make your best SOLVE attempt no later than the second{-}to{-}last turn.} \\
\texttt{Do not wait until the very end.} \\
\\ 
\texttt{You can use these actions:} \\
\texttt{{-} JSON format described at the end of the prompt} \\
\\ 
\\ 
\texttt{Do not add import statements in your guess.} \\
\texttt{Anything you might need is already available, so just write the function body in hypothesis and return the answer.} \\
\\ 
\texttt{Preloaded examples:} \\
\texttt{1. ({-}1, 0)} \\
\texttt{2. (1000, 2)} \\
\texttt{3. (1, 0)} \\
\texttt{4. ({-}2, 2)} \\
\texttt{5. (99, 3)} \\
\texttt{6. ({-}1000, 2)} \\
\texttt{7. ({-}10, 2)} \\
\texttt{8. (2, 2)} \\
\texttt{9. (0, 0)} \\
\texttt{10. (10, 2)} \\
\\ 
\texttt{Scoring:} \\
\texttt{{-} correctness is required for any non{-}zero score} \\
\texttt{{-} in interactive modes, solving with fewer information requests is better} \\
\\ 
\texttt{Input formatting reminders:} \\
\texttt{{-} strings must be quoted} \\
\texttt{{-} booleans must be True/False} \\
\texttt{{-} use commas between multiple values} \\
\texttt{{-} do not write any extra action tags in the same message} \\
\\ 
\texttt{Format example:} \\
\texttt{You write:} \\
\texttt{\{'current\_mode': <probing/confirming/don't know; choose whichever best describes your current state in the task>, 'input\_rationale': <brief description of why you chose this hypothesis>, 'hypothesis': <either the exact string "unknown" if you have no guess yet OR your current best guess as runnable Python matching the given signature, with a one{-}line docstring stating what it does. Never write placeholder code, if you cannot commit to a guess, use "unknown". Match the given signature.} \\
\texttt{e.g.} \\
\texttt{def f(x):} \\
\texttt{	"""Returns the absolute value of x."""} \\
\texttt{	return x if x >= 0 else {-}x} \\
\\ 
\texttt{Use "unknown" if you have no guess.>, 'confidence\_score': <score between 1{-}100. 100 means highest confidence in your hypothesis>\}} \\
\\ 
\texttt{Output only one valid json per turn.} \\
            }
        }
    }
     \\ \\

    \theutterance \stepcounter{utterance}  
    & \multicolumn{4}{p{0.6\linewidth}}{
        \cellcolor[rgb]{0.561,0.737,0.561}{
            \makecell[{{p{\linewidth}}}]{
                \texttt{\tiny{[P1$\rangle$GM]}}
                \texttt{\{"current\_mode":"confirming","input\_rationale":"The outputs match the minimum number of perfect squares needed to sum to |x| for all nontrivial examples: 2=1\^{}2+1\^{}2, 10=3\^{}2+1\^{}2, 1000=30\^{}2+10\^{}2, 99 needs 3 squares. The only anomaly is ±1 mapping to 0, so I special{-}case |x|<=1.","hypothesis":"def f(x):\textbackslash{}n    \textbackslash{}"\textbackslash{}"\textbackslash{}"Returns 0 for |x|<=1, else the minimum number of perfect squares summing to |x|.\textbackslash{}"\textbackslash{}"\textbackslash{}"\textbackslash{}n    n = x if x >= 0 else {-}x\textbackslash{}n    if n <= 1:\textbackslash{}n        return 0\textbackslash{}n\textbackslash{}n    def is\_square(m):\textbackslash{}n        if m < 0:\textbackslash{}n            return False\textbackslash{}n        r = int(m ** 0.5)\textbackslash{}n        while (r + 1) * (r + 1) <= m:\textbackslash{}n            r += 1\textbackslash{}n        while r * r > m:\textbackslash{}n            r {-}= 1\textbackslash{}n        return r * r == m\textbackslash{}n\textbackslash{}n    if is\_square(n):\textbackslash{}n        return 1\textbackslash{}n\textbackslash{}n    a = 1\textbackslash{}n    while a * a <= n:\textbackslash{}n        if is\_square(n {-} a * a):\textbackslash{}n            return 2\textbackslash{}n        a += 1\textbackslash{}n\textbackslash{}n    m = n\textbackslash{}n    while m \% 4 == 0:\textbackslash{}n        m //= 4\textbackslash{}n    if m \% 8 == 7:\textbackslash{}n        return 4\textbackslash{}n    return 3\textbackslash{}n","confidence\_score":63\}} \\
            }
        }
    }
    & & \\ \\

    \theutterance \stepcounter{utterance}  
    & & & \multicolumn{2}{p{0.3\linewidth}}{
        \cellcolor[rgb]{0.9,0.9,0.9}{
            \makecell[{{p{\linewidth}}}]{
                \texttt{\tiny{[GM$|$GM]}}
                \texttt{Game Verdict: LOSS} \\
            }
        }
    }
    & & \\ \\

    \theutterance \stepcounter{utterance}  
    & & & \multicolumn{2}{p{0.3\linewidth}}{
        \cellcolor[rgb]{0.9,0.9,0.9}{
            \makecell[{{p{\linewidth}}}]{
                \texttt{\tiny{[GM$|$GM]}}
                \texttt{Incorrect.} \\
\texttt{Failed Cases:} \\
\texttt{Input: {[}42{]} | Expected: 2 | Got: 3} \\
\texttt{Input: {[}{-}46{]} | Expected: 2 | Got: 3} \\
\texttt{Input: {[}{-}4{]} | Expected: 2 | Got: 1} \\
            }
        }
    }
    & & \\ \\

    \theutterance \stepcounter{utterance}  
    & & & \multicolumn{2}{p{0.3\linewidth}}{
        \cellcolor[rgb]{0.9,0.9,0.9}{
            \makecell[{{p{\linewidth}}}]{
                \texttt{\tiny{[GM$|$GM]}}
                \texttt{def smallest\_divisor\_ge\_two\_abs(x: int) {-}> int:} \\
\texttt{    n = abs(x)} \\
\texttt{    if n < 2:} \\
\texttt{        return 0} \\
\texttt{    for i in range(2, n + 1):} \\
\texttt{        if n \% i == 0:} \\
\texttt{            return i} \\
\texttt{    return 0} \\
            }
        }
    }
    & & \\ \\

    \theutterance \stepcounter{utterance}  
    & & & \multicolumn{2}{p{0.3\linewidth}}{
        \cellcolor[rgb]{0.9,0.9,0.9}{
            \makecell[{{p{\linewidth}}}]{
                \texttt{\tiny{[GM$|$GM]}}
                \texttt{def f(x):} \\
\texttt{    """Returns 0 for |x|<=1, else the minimum number of perfect squares summing to |x|."""} \\
\texttt{    n = x if x >= 0 else {-}x} \\
\texttt{    if n <= 1:} \\
\texttt{        return 0} \\
\\ 
\texttt{    def is\_square(m):} \\
\texttt{        if m < 0:} \\
\texttt{            return False} \\
\texttt{        r = int(m ** 0.5)} \\
\texttt{        while (r + 1) * (r + 1) <= m:} \\
\texttt{            r += 1} \\
\texttt{        while r * r > m:} \\
\texttt{            r {-}= 1} \\
\texttt{        return r * r == m} \\
\\ 
\texttt{    if is\_square(n):} \\
\texttt{        return 1} \\
\\ 
\texttt{    a = 1} \\
\texttt{    while a * a <= n:} \\
\texttt{        if is\_square(n {-} a * a):} \\
\texttt{            return 2} \\
\texttt{        a += 1} \\
\\ 
\texttt{    m = n} \\
\texttt{    while m \% 4 == 0:} \\
\texttt{        m //= 4} \\
\texttt{    if m \% 8 == 7:} \\
\texttt{        return 4} \\
\texttt{    return 3} \\
            }
        }
    }
    & & \\ \\

    \theutterance \stepcounter{utterance}  
    & & & \multicolumn{2}{p{0.3\linewidth}}{
        \cellcolor[rgb]{0.9,0.9,0.9}{
            \makecell[{{p{\linewidth}}}]{
                \texttt{\tiny{[GM$|$GM]}}
                \texttt{Mode: passive\_examples\_oneshot} \\
\texttt{Accuracy: 0.350} \\
\texttt{Efficiency: 1.000} \\
\texttt{Parse errors: 0} \\
\texttt{Runtime errors: 0} \\
            }
        }
    }
    & & \\ \\

\end{supertabular}
}

\end{document}
